\chapter{Cha mẹ và việc học của con}
\markboth{Cha mẹ và việc học của con}{Cha mẹ và việc học của con}

\section{Sao cứ nhắc học hoài?}
\begin{truyen}{Sao cứ nhắc học hoài?}{Family Talk}
\chuhoa{C}{on than: ``Sao ba mẹ cứ nhắc học hoài vậy?''}

Mẹ đáp: ``Vì mẹ sợ con quen với việc buông mình.''

Con cãi: ``Nhắc nhiều chỉ làm con mệt hơn.''

Mẹ im vài giây rồi nói: ``Vậy chắc mẹ cũng phải học lại cách nhắc. Nhưng con cũng phải học cách tự làm để mẹ bớt nhắc.''

\textit{(A child complains that studying is mentioned too often. The parent admits the reminder may be clumsy, but also points out that self-discipline would reduce the need for it.)}
\end{truyen}

\begin{baihoc}
Một nhà đỡ mệt hơn khi người lớn biết nhắc đúng cách và con biết tự giác đúng lúc.
\textit{(A home becomes lighter when adults remind better and children take more responsibility.)}
\end{baihoc}

\begin{ghinhoanh}
\textbf{Short English to remember:}

Too much reminding brings tension.

Self-discipline brings peace.
\end{ghinhoanh}

\begin{tuhocnhanh}
1. Trong nhà em, phần nào nhắc quá nhiều? \textit{(What gets repeated too much at home?)}

2. Con có thể tự làm điều gì để bớt bị nhắc? \textit{(What can the child do alone to reduce reminders?)}
\end{tuhocnhanh}
\ngancach

\section{Điểm thấp có phải con lười?}
\begin{truyen}{Điểm thấp có phải con lười?}{Family Talk}
\chuhoa{S}{au một bài kiểm tra kém, cha hỏi: ``Điểm này là do con lười phải không?''}

Con đáp nhỏ: ``Không hẳn. Có phần con không hiểu thật.''

Cha nói: ``Vậy sao không hỏi sớm?''

Con cúi đầu: ``Tại con sợ bị nói là chậm.''

Cha ngồi yên một lúc mới hiểu rằng nhiều khi con không lười như ông tưởng, chỉ là đã sợ quá lâu.

\textit{(After a low score, a parent assumes laziness. The child reveals a different truth: fear and confusion can hide behind poor results too.)}
\end{truyen}

\begin{baihoc}
Điểm thấp không chỉ có một nguyên nhân. Muốn giúp con, phải hỏi sâu hơn là trách nhanh.
\textit{(Low scores do not have only one cause. To help a child, ask deeper before blaming fast.)}
\end{baihoc}

\begin{ghinhoanh}
\textbf{Short English to remember:}

Low scores may hide fear.

Ask deeper before blaming.
\end{ghinhoanh}

\begin{tuhocnhanh}
1. Em có sợ thú nhận mình không hiểu bài không? \textit{(Are you afraid to admit you do not understand?)}

2. Người lớn trong nhà có hỏi đủ sâu chưa? \textit{(Have the adults at home asked deeply enough?)}
\end{tuhocnhanh}
\ngancach

\section{Con học vì con hay vì ba mẹ?}
\begin{truyen}{Con học vì con hay vì ba mẹ?}{Family Talk}
\chuhoa{C}{on hỏi: ``Rốt cuộc con học vì con hay vì ba mẹ?''}

Ba đáp: ``Lúc đầu có thể vì ba mẹ nhắc. Nhưng đi lâu thì phải là vì con.''

Con nói: ``Nếu con chưa thấy điều đó thì sao?''

Ba trả lời: ``Thì ba mẹ phải giúp con nhìn ra, chứ không chỉ ép con chạy.''

\textit{(A child questions the purpose of learning. The parent admits external pressure may start the journey, but inner ownership must eventually take over.)}
\end{truyen}

\begin{baihoc}
Ép có thể đẩy một người đi vài bước. Chỉ khi hiểu ý nghĩa, người ta mới đi đường dài.
\textit{(Pressure may push someone a few steps. Meaning carries them over the long road.)}
\end{baihoc}

\begin{ghinhoanh}
\textbf{Short English to remember:}

Pressure can start motion.

Meaning sustains the journey.
\end{ghinhoanh}

\begin{tuhocnhanh}
1. Em đang học vì sợ hay vì hiểu? \textit{(Are you studying from fear or understanding?)}

2. Người lớn đã giải thích đủ ý nghĩa chưa? \textit{(Have the adults explained enough meaning?)}
\end{tuhocnhanh}
\ngancach

\section{Học thêm nhiều mới yên tâm?}
\begin{truyen}{Học thêm nhiều mới yên tâm?}{Family Talk}
\chuhoa{M}{ẹ nói: ``Thời nay phải học thêm nhiều mới yên tâm.''}

Con đáp: ``Nhưng con mệt quá.''

Mẹ hỏi: ``Không học thêm thì làm sao cạnh tranh?''

Con nói thật: ``Nếu cạnh tranh mà con kiệt sức thì có đáng không?''

Mẹ im lặng, vì câu hỏi đó không dễ gạt đi.

\textit{(A parent sees extra classes as safety. The child raises the cost: what if the race itself is exhausting the child beyond reason?) }
\end{truyen}

\begin{baihoc}
Sự yên tâm của cha mẹ không nên được đổi bằng sự kiệt sức âm thầm của con.
\textit{(Parental peace of mind should not be bought with a child's quiet exhaustion.)}
\end{baihoc}

\begin{ghinhoanh}
\textbf{Short English to remember:}

More classes are not always more growth.

Exhaustion is a real cost.
\end{ghinhoanh}

\begin{tuhocnhanh}
1. Em có đang quá tải mà không nói ra? \textit{(Are you overloaded without saying it?)}

2. Nhà mình đang theo yên tâm hay theo hiệu quả thật? \textit{(Is your family chasing peace of mind or real effectiveness?)}
\end{tuhocnhanh}
\ngancach

\section{Con mệt mà ba mẹ không thấy?}
\begin{truyen}{Con mệt mà ba mẹ không thấy?}{Family Talk}
\chuhoa{C}{on hỏi trong bữa cơm: ``Ba mẹ có thấy con mệt không?''}

Mẹ khựng lại: ``Mẹ thấy con hay cáu, chứ không biết đó là mệt.''

Con nói: ``Nhiều khi con cáu vì con đuối quá.''

Mẹ thở dài: ``Có lẽ người lớn nhiều lúc chỉ thấy biểu hiện, mà không thấy nỗi mệt phía dưới.''

\textit{(A child asks whether the parents can see the exhaustion underneath irritability. The parent realizes behavior is often noticed before pain is understood.)}
\end{truyen}

\begin{baihoc}
Người thân gần nhau chưa chắc đã nhìn thấy nhau đúng. Thấy biểu hiện chưa phải đã thấy nguyên nhân.
\textit{(Living close does not always mean seeing each other clearly. Seeing the behavior is not the same as seeing the cause.)}
\end{baihoc}

\begin{ghinhoanh}
\textbf{Short English to remember:}

Irritation may hide exhaustion.

Look under the behavior.
\end{ghinhoanh}

\begin{tuhocnhanh}
1. Trong nhà, ai đang bị hiểu sai vì cách họ biểu hiện? \textit{(Who at home is being misunderstood through behavior?)}

2. Em có dám nói thẳng khi mình mệt không? \textit{(Can you say directly when you are tired?)}
\end{tuhocnhanh}
\ngancach

\section{Nếu con không giỏi môn này thì sao?}
\begin{truyen}{Nếu con không giỏi môn này thì sao?}{Family Talk}
\chuhoa{C}{on hỏi: ``Nếu con không giỏi môn này thì sao?''}

Ba đáp hơi gắt: ``Không giỏi thì càng phải cố.''

Con nói: ``Con vẫn cố, nhưng có lẽ con không thể nổi bật ở chỗ này.''

Ba dịu lại: ``Vậy thì cố để không bỏ cuộc, chứ không nhất thiết cố để thành số một.''

\textit{(A child fears not excelling in one subject. The parent slowly shifts from demanding the top position to valuing perseverance.)}
\end{truyen}

\begin{baihoc}
Không phải chỗ nào con cũng phải tỏa sáng. Có chỗ chỉ cần không bỏ cuộc đã là một tiến bộ lớn.
\textit{(A child does not need to shine everywhere. In some places, simply not quitting is already major growth.)}
\end{baihoc}

\begin{ghinhoanh}
\textbf{Short English to remember:}

Not every subject must be your stage.

Persistence is also progress.
\end{ghinhoanh}

\begin{tuhocnhanh}
1. Em đang bị áp lực phải giỏi đều mọi thứ không? \textit{(Are you pressured to excel at everything?)}

2. Mục tiêu hợp lý hơn ở môn yếu là gì? \textit{(What is a more realistic goal in a weak subject?)}
\end{tuhocnhanh}
\ngancach

\section{Cha mẹ kỳ vọng có sai không?}
\begin{truyen}{Cha mẹ kỳ vọng có sai không?}{Family Talk}
\chuhoa{C}{on hỏi: ``Ba mẹ kỳ vọng ở con có sai không?''}

Mẹ nói: ``Kỳ vọng không sai. Chỉ sợ mẹ đặt nặng đến mức làm con nghẹt.''

Con nhìn mẹ: ``Nhiều lúc con thấy vậy thật.''

Mẹ gật đầu: ``Vậy chắc mẹ phải học cách kỳ vọng mà vẫn chừa chỗ cho con thở.''

\textit{(A child questions parental expectation. The parent admits that expectation itself is not wrong, but suffocation is.)}
\end{truyen}

\begin{baihoc}
Yêu thương có định hướng là tốt. Yêu thương biến thành ngột ngạt thì phải dừng lại để sửa.
\textit{(Love with direction is good. Love that turns suffocating must pause and be repaired.)}
\end{baihoc}

\begin{ghinhoanh}
\textbf{Short English to remember:}

Expectation is not always wrong.

Suffocation is.
\end{ghinhoanh}

\begin{tuhocnhanh}
1. Kỳ vọng nào trong nhà đang thành áp lực? \textit{(Which expectation at home has become pressure?)}

2. Làm sao để vẫn có định hướng mà không làm nghẹt nhau? \textit{(How can guidance stay without suffocating one another?)}
\end{tuhocnhanh}
\ngancach

\section{Khi con nói con áp lực}
\begin{truyen}{Khi con nói con áp lực}{Family Talk}
\chuhoa{K}{hi con nói: ``Con áp lực quá,'' ba đáp ngay: ``Có mỗi việc học mà cũng áp lực.''}

Con im luôn.

Một lúc sau mẹ nói nhỏ với ba: ``Có thể với mình là bình thường, nhưng với con thì đang quá sức.''

Ba nhìn sang con và hiểu ra rằng so nỗi mệt của mình với nỗi mệt của con không làm nỗi mệt nào nhẹ đi cả.

\textit{(A child says the pressure feels too much. An adult dismisses it at first, then realizes comparison does not reduce another person's burden.)}
\end{truyen}

\begin{baihoc}
Muốn con nói thật, người lớn phải bớt coi nhẹ cảm xúc mà mình không trực tiếp mang.
\textit{(If adults want honesty from children, they must stop minimizing feelings they do not personally carry.)}
\end{baihoc}

\begin{ghinhoanh}
\textbf{Short English to remember:}

Do not compare pain to dismiss it.

Listening keeps honesty alive.
\end{ghinhoanh}

\begin{tuhocnhanh}
1. Em từng bị coi nhẹ áp lực chưa? \textit{(Have your pressures ever been dismissed?)}

2. Trong nhà, ai cần học cách nghe hơn? \textit{(Who at home needs to learn listening more?)}
\end{tuhocnhanh}
\ngancach

\section{Ba mẹ chỉ hỏi điểm thôi}
\begin{truyen}{Ba mẹ chỉ hỏi điểm thôi}{Family Talk}
\chuhoa{C}{on nói một câu rất thật: ``Mỗi lần đi học về, ba mẹ chỉ hỏi điểm thôi.''}

Ba ngạc nhiên: ``Chứ còn hỏi gì nữa?''

Con đáp: ``Hỏi hôm nay con có ổn không. Hỏi con có chuyện gì vui không. Hỏi con đang lo gì cũng được.''

Ba im rất lâu, vì nhận ra mình đã yêu con bằng cách quá hẹp.

\textit{(A child feels reduced to grades. The parent suddenly sees how narrow care can become when it only asks for measurable results.)}
\end{truyen}

\begin{baihoc}
Khi chỉ hỏi thành tích, người lớn rất dễ làm con tưởng mình chỉ đáng giá ở thành tích.
\textit{(When adults ask only about performance, children may think their worth lives only there.)}
\end{baihoc}

\begin{ghinhoanh}
\textbf{Short English to remember:}

Children need more than score questions.

Love must ask wider questions.
\end{ghinhoanh}

\begin{tuhocnhanh}
1. Em muốn được hỏi thêm điều gì ngoài điểm số? \textit{(What else do you want to be asked beyond grades?)}

2. Người lớn trong nhà đang hỏi quá hẹp ở đâu? \textit{(Where are the adults asking too narrowly?)}
\end{tuhocnhanh}
\ngancach

\section{Có phải học giỏi mới làm cha mẹ yên lòng?}
\begin{truyen}{Có phải học giỏi mới làm cha mẹ yên lòng?}{Family Talk}
\chuhoa{C}{on hỏi: ``Có phải con học giỏi thì ba mẹ mới yên lòng không?''}

Mẹ đáp thật: ``Mẹ yên lòng hơn khi con biết sống tử tế và không bỏ mình. Học giỏi thì tốt, nhưng không phải tất cả.''

Con nhìn mẹ như để kiểm tra xem câu đó có thật không.

Mẹ mỉm cười: ``Mẹ cũng đang học cách nhớ điều đó mỗi ngày.''

\textit{(A child wonders if high achievement is the only path to parental peace. The parent answers with a fuller standard: character and self-respect matter more.)}
\end{truyen}

\begin{baihoc}
Thành tích nên là một phần của niềm tự hào, không nên là toàn bộ điều kiện để được yêu thương.
\textit{(Achievement should be one part of pride, not the full condition for being loved.)}
\end{baihoc}

\begin{ghinhoanh}
\textbf{Short English to remember:}

Grades matter, but not alone.

Character matters more in the long run.
\end{ghinhoanh}

\begin{tuhocnhanh}
1. Em đang hiểu tình thương của cha mẹ theo điều kiện nào? \textit{(What conditions do you think your parents place on love?)}

2. Trong nhà, điều gì cần được nói rõ hơn để bớt hiểu lầm? \textit{(What should be said more clearly at home?)}
\end{tuhocnhanh}
\ngancach